% !TeX root = ../../main.tex

\subsection{英語の分析}\vspace{-0.2em}
英語の入試について、分析していく。適宜、別冊付録の分析結果を参照してほしい。

問題構成は大問1から4までの4つの大問構成である。

大問1はリスニング。配点は35点。

大問2は語句・文法問題。配点は16点。

大問3はさらにAからCまでの3つの小問に分かれており、それぞれに長文読解問題がある。それぞれの小問の配点は10点、10点、17点と
やや3Cの配点が大きい。

大問4は英作文問題。配点は12点。
\vspace{1em}

分析していて感じたことは、リスニングの配点が大きいことはもちろんのこと、
\textbf{空欄補充及び英作文の問題数が大きと感じた。}英語といえば、「長文読解」というイメージだが、
これらの問題数が多いことは、英語の基本を問う問題が多いことを示している。英語が苦手な生徒と得意な生徒の差がつきにくいのではないかと思う。

一方、長文読解は文量はそこまで多くないにしろ、得手不得手がはっきりと出る問題であるため、差がつきやすい。

\vfill
